\chapter{Conclusion}

A central motivation of this work was the question of whether the
ability of entangled quantum states to represent non-separable
dependencies can provide advantages when modeling the correlations
between multiple underlyings on which basket options depend.
Against this background, this thesis investigated the extent to
which VQCs are suitable for approximating basket option prices.
To assess their performance, a comparison with established classical
neural network architectures was conducted.

The results show that both classical neural networks and VQC-based
models are able to approximate the underlying pricing function of
basket options with high accuracy. In the baseline scenario with a
large training dataset and without artificial noise, almost all
models achieve test-$R^2$ values close to $1$. This confirms that
data-driven approaches are generally well suited to learn complex
valuation functions in the context of option pricing.

The analysis of model capacity reveals notable differences between
classical architectures and VQC models. While classical neural
networks tend to show continuous performance improvements as the
model size increases, VQCs already achieve most of their maximum
performance at moderate circuit depth. Additional layers expand
the representable frequency spectrum, but provide only limited
benefit in the considered datasets since the target functions are
dominated by low-frequency components.

Under limited data availability, a clear trade-off between model
capacity and robustness can be observed. Simpler architectures --
both classical networks and VQCs -- produce more stable results for
small datasets, whereas more complex models reach their full
potential only when sufficient training data are available.
A similar pattern appears under noisy training data, where models
with higher capacity suffer stronger performance degradation.
This effect becomes particularly pronounced in the combined
scenario of limited data and strong noise, where simpler models,
especially shallow VQCs and smaller classical networks, show the
most robust performance.

The analysis of a temporally ordered train–test split further
reveals differences in generalization behavior. While classical
ReLU-based networks largely retain their predictive performance,
VQCs show greater sensitivity to distribution shifts in the
input features, particularly at higher volatility levels.
This suggests that the globally coupled functional structure
of rotation-based quantum models allows highly accurate
approximation within the training distribution but may lead
to less robust extrapolation.

Overall, the results demonstrate that VQC-based models can achieve
competitive approximation performance for basket option pricing.
However, no clear practical advantage over classical neural
networks can be identified in the investigated setting.
At the same time, the experiments highlight that their performance
strongly depends on the structure of the target function as well
as on data availability, data quality, and the stability of the
data distribution.